\subsection{Comparaison des assemblages chromosomiques entiers}
Les versions v1.0 et v2.0 de CHM13 permettent de mesurer le coût d'une révision d'assemblage sur le chromosome 20. Les entrées \texttt{CP068258.1} et \texttt{CP068258.2} comptent 66\,210\,247 et 66\,210\,255 bases~\cite{chmreleases}. Le chromosome 20 porte déjà l'accession \texttt{CP068258.2} dans la version intermédiaire 1.1. Sur le chromosome 21, GRCh38 et CHM13 représentent deux références distinctes et comportent des bases ambiguës traitées comme des caractères ordinaires.

Une borne exploite l'histogramme des lettres, où \(f_A(c)\) désigne l'effectif de la lettre \(c\) dans \(A\). Pour \(u_c=f_A(c)-f_B(c)\) et \(x_+=\max(0,x)\), chaque édition efface au plus une unité d'excédent et une de déficit. Ainsi,
\[
 d(A,B)\ge\max\!\left(\sum_c(u_c)_+,\ \sum_c(-u_c)_+\right).
\]
Une borne supérieure suit un alignement global réalisable. Sur chaque bloc de la chaîne UCSC 23, il paie les désaccords. Dans chaque intervalle entre blocs, il choisit le moins coûteux des deux alignements ancrés à gauche ou à droite, puis paie l'écart de longueur. Les extrémités suivent la même règle. La somme des coûts des blocs et des intervalles donne le coût du chemin entier.
\chart{biology/chromosomes_entiers.png}{Les durées sur le chromosome 20 et les bornes de distance sur le chromosome 21 comparent deux niveaux de divergence entre assemblages.}{fig:chrentiers}{bench/human/whole_chromosomes.py}
Le panneau des temps de la figure~\ref{fig:chrentiers} porte sur cinq blocs à chaud sur un cœur Grace. Les deux bibliothèques donnent \(d=30\) sur le chromosome 20. Edlib prend \(0{,}9767\pm0{,}0026\)\,s et WFA2 \(0{,}01644\pm0{,}00012\)\,s. Pour le chromosome 21, les appels de la campagne EPYC ont tous atteint la limite murale de 120\,s. L'écart de longueur seul vaut 1\,619\,301 bases.

Sur le chromosome 20, un chemin optimal donne \(N_{\mathrm{sub}}=18\), \(N_{\mathrm{ins}}=10\) et \(N_{\mathrm{del}}=2\). Leur somme vaut 30, tandis que les huit bases de différence de longueur correspondent à \(10-2\). Sur ce chemin, 19 éditions précèdent un bloc identique de 66\,209\,720 bases et 11 le suivent. Les intervalles de base zéro, fin exclue, sont \(A[270,66\,209\,990)\) et \(B[277,66\,209\,997)\). Le bloc couvre 99,999204\,\% de la première version. Cette localisation décrit une révision d'assemblage et l'un des chemins optimaux possibles. Le consortium T2T signale une amélioration du polissage des régions télomériques entre les versions 1.0 et 1.1~\cite{chmreleases,mccartney}. Les coordonnées du chemin ne suffisent pas pour attribuer chaque édition à une correction particulière. La figure~\ref{fig:chrentiers} situe les opérations par rapport au bloc commun.

\synthese{Pour le chromosome 21 entier, la composition et un chemin réalisable encadrent la distance entre 6\,621\,364 et 16\,015\,052 éditions. Le score exact reste inconnu après les arrêts à l'échéance.}
